\begin{center}\begin{tikzpicture}[font=\small,>=stealth,line join=round]
\draw[gray!40,rounded corners=3pt] (0,0) rectangle (14.2,8.3);\node at (7.1,7.700000000000001) {左右等距、函数值同高：对称轴就在中间};
\begin{scope}[shift={(1.7923076923076926,2.038732394366197)},x=1.153846153846154cm,y=0.647887323943662cm]\draw[->,gray] (-0.6,0)--(4.6,0) node[right]{$x$};\draw[->,gray] (0,-0.6)--(0,6.5) node[above]{$y$};\draw[blue,thick,solid,domain=-0.1:4.1,samples=180] plot(\x,{(\x-2)^2+1});\fill[orange!85!black] (0,5) circle(2.2pt);\fill[orange!85!black] (1,2) circle(2.2pt);\fill[orange!85!black] (2,1) circle(2.2pt);\fill[orange!85!black] (3,2) circle(2.2pt);\fill[orange!85!black] (4,5) circle(2.2pt);\draw[violet,dashed] (2,-.5)--(2,6);\node[above,violet] at (2,6){$x=2$};\draw[orange!85!black,dashed] (0,5)--(4,5);\draw[orange!85!black,dashed] (1,2)--(3,2);\foreach \v in {0,1,2,3,4}{\draw[gray] (\v,0)--(\v,-.1);\node[below=5pt] at (\v,-.1){$\v$};}\end{scope}\node[anchor=west,align=left,text width=5.0cm] at (8.4,4.0) {$f=(x-2)^2+1$\\[9pt]$f(0)=f(4)=5$\\[9pt]$f(1)=f(3)=2$\\[9pt]每对横坐标之和为 $4$\\[9pt]轴的位置：$4\div2=2$};
\node[text width=12.6cm,align=center] at (7.1,.65) {图示具体配对；完整证明要求定义域左右成对，且每一对都同高。};
\end{tikzpicture}\end{center}
